% Diagrama de máquina de estados (UML 2) do caso. Elementos e fontes no código:
% docs/sistema/diagramas-fontes.md, seção 8.
\begin{figure}[htbp]
\centering
\caption{Diagrama de máquina de estados do caso, com o que cada estado permite}
\label{fig:uml-estado-caso}
\resizebox{\textwidth}{!}{%
\begin{tikzpicture}[
  font=\footnotesize,
  estado/.style={draw, rounded corners=6pt, rectangle split, rectangle split parts=2,
    align=left, text width=#1, inner sep=4pt, fill=white},
  inicial/.style={circle, fill=black, inner sep=0pt, minimum size=7pt},
  trans/.style={-{Straight Barb[length=6pt]}},
  rot/.style={font=\scriptsize, align=center, inner sep=1.5pt, fill=white},
  nota/.style={align=left, font=\scriptsize, inner sep=4pt}
]
\newcommand{\dobraest}[1]{%
  \draw (#1.north west) -- ([xshift=-6pt]#1.north east) -- ([yshift=-6pt]#1.north east)
        -- (#1.south east) -- (#1.south west) -- cycle;
  \draw ([xshift=-6pt]#1.north east) |- ([yshift=-6pt]#1.north east);}
\node[estado=6.2cm, anchor=north] (ativo) at (0,0) {%
  \centering\textbf{Ativo}\par{\scriptsize hoje $\leq$ \texttt{retencao\_ate}\par ou caso sem prazo\par}
  \nodepart{second}\raggedright\scriptsize
  \textbf{Permite}, além do que Somente leitura permite:\\
  -- coletar e recoletar;\\
  -- buscar entidades, ver detalhe, grafo e caminho;\\
  -- exportar o grafo (GraphML, CSV);\\
  -- executar tipologias e decidir alertas;\\
  -- ler o texto de captura, extrair com IA e decidir sugestões;\\
  -- anotar; criar e alterar tarefas;\\
  -- criar, reativar e executar monitoramentos;\\
  -- o agendador executa os monitoramentos ativos.};
\node[estado=7.4cm, anchor=north] (leit) at (10.6,0) {%
  \centering\textbf{Somente leitura}\par{\scriptsize hoje $>$ \texttt{retencao\_ate}\par (prazo de retenção vencido)\par}
  \nodepart{second}\raggedright\scriptsize
  \textbf{Permite} (sem \texttt{caso\_aberto}):\\
  -- ver o caso; listar coletas, comparações, alertas, membros,
  anotações, tarefas, avisos, monitoramentos e sugestões;\\
  -- ver, comparar e verificar capturas; gerar o relatório;\\
  -- gerenciar membros (papel responsável);\\
  -- desativar monitoramento; marcar aviso como lido.\\[3pt]
  \textbf{Recusa} com 403 e \texttt{caso.retencao\_recusa} na trilha:\\
  -- tudo o que só Ativo permite.\\[3pt]
  \textbf{Agendador}: execução \texttt{recusada\_retencao}, monitoramento
  desativado e \texttt{monitoramento.retencao\_recusa} na trilha.};
\node[inicial] (i0) at ($(ativo.north)+(0,1.3)$) {};
\draw[trans] (i0) -- node[rot, right=2pt, align=left]{criar caso (\texttt{finalidade}, \texttt{retencao\_ate}) /\\
  criador vira responsável; \texttt{caso.criar} na trilha} (ativo.north);
\draw[trans] (ativo.east |- leit.text split) -- node[rot, above, align=center]{when(hoje UTC $>$\\\texttt{retencao\_ate})}
  (leit.west |- leit.text split);
\node[nota, text width=16.0cm, anchor=north west] (nt) at ($(ativo.south west |- leit.south)+(0,-0.35)$) {%
  Nos dois estados vale também o papel no caso: leitura exige leitor; alteração, investigador;
  membros, responsável. O prazo (\texttt{retencao\_ate}) só é gravado na criação e o último dia
  ainda vale; por isso não há volta a Ativo. Não há estado de exclusão: nada é apagado pelo
  sistema, e o descarte depois do prazo fica a cargo do órgão responsável.};
\dobraest{nt}
\end{tikzpicture}}
\fonte{Elaborada pelo autor a partir de \texttt{buscador/retencao.py}, \texttt{buscador/api/deps.py}
(\texttt{Contexto.caso\_aberto}), das rotas em \texttt{buscador/api/rotas\_*.py} (quais chamam
\texttt{caso\_aberto} e quais só \texttt{Contexto.caso}) e de \texttt{buscador/api/monitoramento.py}.}
\end{figure}
